\section{Agreement of local steps}\label{sec:local-comparison}

\subsection{Common lifts and local algebras}
\label{subsec:rgcov-common-lifts}

For the rest of the comparison assume
\begin{equation}\label{eq:rgcov-comparison-side}
                      \ell_1>(31L+16)M.
\end{equation}
Put $d_U=2L+1$, $c_l=2L$.
For a finite nonempty union $\widetilde X$ of covering
$M$-cubes, let $\operatorname{box}(\widetilde X)$ be
its smallest rectangular union of cubes and let
$\operatorname{diam}(\widetilde X)$ be that box's
largest side. Define
\[
 N(\widetilde X)=\operatorname{box}(\widetilde X)^{\sim6d_U}.
\]
We call $\widetilde X$ admissible
when $\operatorname{diam}(\widetilde X)
\le\ell_1-14d_UM$; a polymer is flat when it is the
image of a connected admissible $\widetilde X$.
These are the sets on which the two torus algebras
will be identified.

\begin{proposition}\label{prop:rgcov-common-lifts}
For admissible $\widetilde X$, projection is injective
on $N(\widetilde X)$ and on every supported index
contained there. An index of either torus supported
in the image of a subset of $N(\widetilde X)$ has
exactly one lift contained in that subset.
Pairs have their shortest lifts, with the same
$d_X$ and six paths.
Local incidences of total support diameter below
$2d_UM$ are preserved in both directions.

There is a norm-preserving algebra isomorphism
$\iota_{\widetilde X}$ between the two subalgebras
generated by ordinary fields over $N(\widetilde X)$,
pair fields with both endpoints there, and their
unit fluctuation and auxiliary fields there.
It identifies covering generator indices.
At zero it includes the chamber $W$ fields and
no auxiliary $\Psi$.
\end{proposition}

\begin{proof}
Every side of $N(\widetilde X)$ is at most
$\ell_1-2d_UM$. Since faces contain no fine lattice
sites, any two of its sites satisfy
\[
 |u-u'|_\infty\le\ell_1-2d_UM-\eta_k.
\]
For $w\ne0$ this implies
\begin{equation}\label{eq:rgcov-lift-separation}
              |u-u'-\ell_iw|_\infty\ge2d_UM+\eta_k.
\end{equation}
It proves disjointness of $N(\widetilde X)$ from
its nonzero $\ell_1$ translates, injectivity, and
uniqueness of any representative with distance
below $2d_UM$.

To lift an index supported in the projected subset,
translate one of its lifts so that one support point
lies in that subset. Every remaining support point
is congruent to a point of the subset and is within
$2LM<2d_UM$ of the first. Equation~
\eqref{eq:rgcov-lift-separation} forces the congruence
translation for each point to be zero.
Thus the whole index lies in the subset; injectivity
gives uniqueness. In particular the shortest pair
has its endpoints there, distance at most $M$, and
its same covering paths.

For an incidence relation, lift a tuple realizing
the projected relation and translate its first index
to the prescribed lift. Each other lifted index
differs from the prescribed one by an $\ell_i$
period. Its support is at distance below $2d_UM$
from the first one. Equation~
\eqref{eq:rgcov-lift-separation} again forces the
period to vanish. This proves the reverse incidence
implication; the forward one is projection.
The relations used below meet this range condition:
nearest neighbors and chamber neighbors, short
pairs and their paths, block and cube membership,
common faces, linked parameter cubes, cube meeting,
and coordinate support. The largest diameter is
$4LM<2d_UM$, for two linked photon parameter cubes.
They are all invariant under $LM\Z^3$ translations.

No two covering generator indices in $N(\widetilde X)$
have the same period class on either torus.
Relative to any fundamental generator system their
images are therefore independent generators, up
to their antiperiodic signs.
Identifying the same covering indices extends
uniquely to a free Grassmann algebra isomorphism.
It preserves bar indices and degrees and sends
basis monomials to basis monomials, up to the
same coefficient signs. All generator weights and
index measures are common, so it preserves norms.
The maps for nested admissible sets restrict to
one another.
\end{proof}

A local activity on $X_i=\operatorname{pr}_i\widetilde X$
belongs to this identified algebra. Ordinary sites
lift to $\widetilde X$; an admissible pair lifts
to $\widetilde X^{\sim1}$. For its first kind the
two sites and paths lie in $\widetilde X$.
For its second kind the pair is $(x,x^{(3)})$
with an endpoint over $\widetilde X$.
Each assertion follows from the incidence comparison.
Two local families $F_1,F_2$ agree on flat polymers
when
\begin{equation}\label{eq:rgcov-local-agreement}
 \iota_{\widetilde X}F_1(X_1;A_1,\mathscr C_1)
                             =F_2(X_2;A_2,\mathscr C_2)
\end{equation}
whenever the two admissible connection pairs have
equal lifts on $N(\widetilde X)$.
It suffices to check $A_2=p^*A_1$,
$\mathscr C_2=p^*\mathscr C_1$:
for other pairs with equal lifts, locality of
$F_2$ identifies them with these pullbacks on all
bonds meeting $\widetilde X$.
Thus once agreement is proved, equality holds
already under equality on those read bonds.

\begin{theorem}\label{thm:rgcov-step-comparison}
Suppose the two sides at levels $k$ and $k+1$ are
admissible and satisfy
$\ell_1/L>(31L+16)M$. Use the projected and top
systems constructed above, with the hypotheses of
Theorem~\extref{A}{NonlinearStep} on both tori.
Let the real mass $m$ be common with $|m|\le w_k^{\mathrm m}$,
and let the incoming real activities satisfy
$\|\mathscr E_i\|_k^x\le w_k^{\mathrm E}$ and agree on flat polymers,
with $\mathscr E_i=0$ at zero. For the external scalar conclusion,
assume the incoming scalar families agree on flat polymers as well.
Then their extracted mass and vacuum contributions
are equal. Their outgoing activities agree on flat
polymers at level $k+1$.
The external scalar recursion
\eqref{eq:rgnl-external-scalar-map} preserves agreement
of scalar families on such polymers.
All these assertions hold for the activities without
an external slot as well.
\end{theorem}

The proof is given in \cref{subsec:rgcov-step-comparison}, after
comparison of the local operations.

Every small polymer is flat. The net diameter bound
gives
$\operatorname{diam}(\widetilde X)
\le M(d_M(\widetilde X)+3)\le M(L+3)$ for a small
connected lift. By \eqref{eq:rgcov-comparison-side},
$M(L+3)+14d_UM<\ell_1$.
The nonwrapping lift proposition then identifies
small torus polymers containing a chosen cube
with small covering polymers containing its
chosen lift. This identification is the same
on both tori. Conversely, if an admissible
connected lift projects to a small polymer,
the small nonwrapping lift is its period translate:
start at one common cube and propagate equality
along adjacent cubes using the incidence result.
Hence smallness, tree length and volume coincide
for its two images.

Subpolymers and parameter cubes contained in an
admissible $\widetilde Y$ correspond bijectively
on the two tori. Common faces and links are
preserved, so incompatibility is preserved.
The definitions by cube meeting give
\[
 P_0(\operatorname{pr}_i\widetilde X)
                 =\operatorname{pr}_iP_0(\widetilde X),\qquad
 \operatorname{pr}_iY'(\widetilde X,\widetilde P)
          =Y'(\operatorname{pr}_i\widetilde X,\operatorname{pr}_i\widetilde P).
\]
When $\widetilde X$ and all cubes of $\widetilde P$
are contained in $\widetilde Y$, the resulting
$Y'$ is contained there too, since the $M$-cubes
partition sites. Reblocking satisfies
$\overline{\widetilde X}^{\,L}
\subset\widetilde Y^{\sim(L-1)}$ and commutes with
projection. These larger sets still lie in
$N(\widetilde Y)$, so their incidences have the
same unique lifts.

\subsection{Supported rows near a flat set}
\label{subsec:rgcov-supported-rows}

The preceding identification applies to the weakened
primitives because setting parameters to zero separates
their components. We first record the component argument.
Let $s$ vanish at every parameter cube not contained
in a union $Z$ of $M$-cubes. If $Z_1$ is a face-connected
component of $Z$ and
$\square(t)^{\sim2L}\subset Z_1$, then
\begin{equation}\label{eq:rgcov-component-row}
 \begin{split}
 O(s,A)(t,t')&=0 \quad\text{unless }t'\text{ is contained in }Z_1,\\
 O(s,A)(t,\cdot)&=O(s\mathbf1_{Z_1},A')(t,\cdot).
 \end{split}
\end{equation}
Here $O$ is a parameter primitive, and $A'$ is any
connection in the primitive's domain equal to $A$ on bonds meeting $Z_1$;
the photon primitives have no connection argument.
The set $\square(t)$ is the union of the $M$-cubes
meeting the support of $t$.

To prove \eqref{eq:rgcov-component-row}, choose a parameter
cube containing that support inside
$\square(t)^{\sim2L}$. For a chamber site or a unit site
it is its $M$-cube. For a fine bond it is the parameter
cube of radius $M$ about the $M$-cube center of one end:
both ends lie in it, and it is contained in
$\square(t)^{\sim1}$. For a photon coordinate index choose
the parameter cube of radius $LM$ about the $LM$-cube
center of one point of its support.
The support diameter is at most $L$ and $M\ge2$,
so this cube contains the entire support.
If $z_c$ is the $M$-cube center of that point and $z''$
is the $M$-cube center of any site in the parameter cube,
then
\[
 |z''-z_c|_\infty
 <(3L/2+1)M\le(2L+1)M.
\]
These centers differ by a multiple of $M$, hence the
cube is contained in $\square(t)^{\sim2L}$.

Let $\mathscr Y$ be all parameter cubes contained in $Z$,
and let $\mathscr Y_1$ be their link component containing
the cube just chosen. The $M$-cubes meeting
$\bigcup\mathscr Y_1$ form a face-connected polymer
contained in $Z$. It meets $\square(t)\subset Z_1$,
and therefore is contained in the component $Z_1$.
The component restriction property of the weakened
primitives, proved in the photon and fermion constructions,
now confines the row to $\bigcup\mathscr Y_1$.
Replacing $s$ by $s\mathbf1_{Z_1}$ leaves all its
parameters there unchanged; replacing $A$ by $A'$
leaves all its read bonds unchanged. This proves
\eqref{eq:rgcov-component-row}.

\begin{proposition}\label{prop:rgcov-supported-rows}
Let $\widetilde Y$ be admissible, $Y_i=\operatorname{pr}_i
\widetilde Y$, and $0\le j\le5d_U$. Let $s_1$ vanish
outside $Y_1$. Define $s_2$ by giving the corresponding
parameter cubes contained in $\widetilde Y$ the same
values and setting all other entries to zero.
Suppose $A_i\in\widetilde{\mathcal G}_k(T_i)$ have equal
lifts on the bonds contained in
$\widetilde Y^{\sim(j+d_U)}$.
For each parameter primitive and each covering row $t$
contained in $\widetilde Y^{\sim j}$, the two rows,
written with covering indices in $N(\widetilde Y)$,
coincide:
\begin{equation}\label{eq:rgcov-local-primitive-row}
                 O_1(s_1,A_1)(t,t')=O_2(s_2,A_2)(t,t').
\end{equation}
Both rows have columns contained in
$Y_i^{\sim(j+c_l)}$, and read the connection only on
bonds contained in $Y_i^{\sim(j+d_U)}$.

For a connected nonempty $\widetilde X\subset\widetilde Y$
and connections with equal lifts on bonds meeting
$\widetilde X$, the local polymer families also satisfy
\begin{equation}\label{eq:rgcov-local-polymer-kernel}
 O_1(\operatorname{pr}_1\widetilde X,A_1)(t,t')
       =O_2(\operatorname{pr}_2\widetilde X,A_2)(t,t')
\end{equation}
at rows and columns over $\widetilde X$.
Here $O=M_k$ for $k\ge1$, or $O=S_{k,y}$ for any $k$
and $y\ge0$.
\end{proposition}

\begin{proof}
Put $\widetilde B=\widetilde Y^{\sim(j+c_l)}$.
Its one-cube enlargement is
$\widetilde Y^{\sim(j+d_U)}\subset N(\widetilde Y)$.
For the row $t$, the set $\square(t)$ is one $M$-cube
or two face-adjacent $M$-cubes. This follows directly
for site and bond indices; a coordinate support lies
in one $L$-block or is the support of one bond.
Thus $\square(t)^{\sim c_l}$ is a connected box
contained in $\widetilde B$.
Let $\widetilde B_1$ be the component of $\widetilde B$
containing this box.

Its image is a component of
$\operatorname{pr}_1\widetilde B$.
On $T_2$ its image is a component both of
$\operatorname{pr}_2\widetilde B$ and of
$p^{-1}(\operatorname{pr}_1\widetilde B)$.
Indeed a cube adjacent to its image has a lift adjacent
to a cube of $\widetilde B_1$. This lift lies in
$\widetilde B^{\sim1}\subset N(\widetilde Y)$.
If it belongs to a period translate of $\widetilde B$,
separation \eqref{eq:rgcov-lift-separation} forces that
translate to be zero. It then lies in $\widetilde B$
and in the same component. This proves maximality
of the images as components.

Apply \eqref{eq:rgcov-component-row} on $T_1$.
It confines the row to $\operatorname{pr}_1\widetilde B_1$
and lets the connection be changed off the bonds meeting
that component. Such bonds lift into
$\widetilde B^{\sim1}$, proving the claimed read set.
On the parameter cubes contained in
$\operatorname{pr}_2\widetilde B_1$, the families
$p^*s_1$ and $s_2$ agree. To see this even at cubes
not contained in $\widetilde Y$, observe that each
such cube in $N(\widetilde Y)$ has a unique lift;
its image cannot also be the image of a cube in
$\widetilde Y$. Both entries are therefore zero.
The connections $p^*A_1$ and $A_2$ agree on bonds
meeting $\operatorname{pr}_2\widetilde B_1$.
Applying \eqref{eq:rgcov-component-row} twice on $T_2$
gives
\[
 O_2(p^*s_1,p^*A_1)(t,\cdot)
       =O_2(s_2,A_2)(t,\cdot),
\]
with columns confined to $\operatorname{pr}_2\widetilde B_1$.

In the signed image formula
\eqref{eq:rgcov-signed-image-kernel}, a nonzero summand
at $t'+\ell_1w$ would have a lift in $\widetilde B_1$.
Both this lift and $t'$ belong to $N(\widetilde Y)$;
their difference is an $\ell_1$ period.
Separation forces $w=0$ in $(\Z/J_c\Z)^3$.
That summand has character $1$, proving
\eqref{eq:rgcov-local-primitive-row}.

For the polymer families, every polymer $X'$ on $T_2$
with $pX'=\operatorname{pr}_1\widetilde X$ is
$\operatorname{pr}_2(\widetilde X+v)$ for one
$v\in\ell_1\Z^3/\ell_2\Z^3$.
Distinct classes give disjoint cube sets.
To verify the assertion for an arbitrary $X'$, lift
two adjacent cubes as cubes of $\widetilde X+v_1$
and $\widetilde X+v_2$. Their adjacency places a
translate of the second cube in
$\widetilde X^{\sim1}\subset N(\widetilde Y)$.
Separation forces $v_1=v_2$ modulo $\ell_2\Z^3$.
Connectedness gives one class for every cube of $X'$,
and projection onto the whole image gives equality
with that entire copy.

The projected kernel is consequently a double sum over
$v$ and the column image $w$. The local polymer support
requires both its row and column to be over
$\operatorname{pr}_2(\widetilde X+v)$.
The common row lift forces $v=0$; the common column
lift then forces $w=0$. The remaining term is
$O_2(\operatorname{pr}_2\widetilde X,p^*A_1)(t,t')$.
Locality on the bonds meeting this polymer permits
$p^*A_1$ to be replaced by $A_2$.

The barred mean field satisfies
$\bar H_k(s,A)=\mathsf C^{-1}H_k(s,-A)\mathsf C$,
so its comparison follows at the connections $-A_i$.
At zero the mean-field kernels are
$48\delta_{x,x'}\mathbf1_4$ times their parameter products;
the coordinate and chamber measures are the same
on the two tori, and the chamber label is preserved.
There is no $M_0$ family. The proof applies to
$\Gamma_0$ and $S_{0,y}$ on the full chamber indices.
\end{proof}

\subsection{Finite rows and the local operations}
\label{subsec:rgcov-finite-rows}

For finite operators, equality of rows follows directly
from their covering formulas. Suppose a row has all its
columns in $N(\widetilde Y)$, and its connection and
polymer membership data agree on the bonds and sites
which that row reads. The two formulas are then the
same finite sum on either torus. There are no repeated
column classes by \cref{prop:rgcov-common-lifts}.
Writing the row as
$\sum_{t'}r(t,t')f(t')$, its kernel is
\begin{equation}\label{eq:rgcov-row-kernel-weight}
                           O(t,t')=\nu(t')^{-1}r(t,t').
\end{equation}
Thus equality holds for the kernels as well, with the
same column weight.

We need this observation at connections $B_i$ having
equal lifts on bonds meeting a connected
$\widetilde U\subset\widetilde Y$.
The photon coordinate map $C$ reads no connection:
a noncentral face bond is its own coordinate, and
the central bond is the negative sum of the other
face bonds and the interior straight-path contributions
$n_c(b)/L$. All columns are supported in the one or
two $L$-blocks of the row bond's ends. Hence its rows
at bonds in $\widetilde U^{\sim(1+2d_U)}$ agree and
have range at most $2L-1$.

The covariant difference row at an admissible pair is
\[
 d_X(x,x')^{-\alpha}
               \big(U_B(x,x')f(x')-f(x)\big).
\]
Its only columns are $x,x'$, and it reads exactly
the six pair paths. These paths are contained in
$\widetilde U$ for the first kind, or are the bond
meeting $\widetilde U$ for the second kind.
The range is $M$; inside a chamber cluster the
transport is $1$ and no bond is read.
These rows therefore agree, also at $-B_i$.

For $k\ge1$ the reconstruction rows $T^{(1)},T^{(2)}$
and their barred versions are the two-piece formulas
of the reconstruction estimate. At $y\in\widetilde U$
they have columns over $\Delta_y$ and at $x^{(3)}$,
or at the pair $(x,x^{(3)})$ across its boundary.
They read the bonds in $\Delta_y$ and the bonds
$(c,c_3)$ starting there. All meet $\widetilde U$.
Their membership indicators are the same because
\[
 u\in\widetilde U+\ell_i\Z^3
       \quad\Longleftrightarrow\quad u\in\widetilde U
             \qquad(u\in\widetilde U^{\sim1}),
\]
again by separation. Every site involved is at
distance at most $(1+\eta_k)/2<1$ from $y$.
This proves equality of both pieces, including
the boundary term.

The finite rows of $Q_kJ^*$ and $JQ_k^T(-B)$ have
their columns and read bonds in a single $\Delta_y$.
The row of $B_{k,r}$, where $r\ge0$ is the resolvent
parameter, has its columns in one $L$-block and reads
the paths from its center to its sites.
These paths lie in that block. Their kernels are
\begin{align*}
 (Q_kJ^*)(y,x)&=\mathbf1[c\in\Delta_y]\varphi_B(y,c)\mathbf1_4,\\
 (JQ_k^T(-B))(x,y)&=\mathbf1[c\in\Delta_y]\varphi_B(y,c)^{-1}\mathbf1_4,\\
 B_{k,r}
   &=(b_k+\mathrm i\gamma_5r)^{-1}I\\
   &\quad+\big((b_k+b_L+\mathrm i\gamma_5r)^{-1}
                    -(b_k+\mathrm i\gamma_5r)^{-1}\big)P_B .
\end{align*}
The projection $P_B$ has kernel $L^{-3}$ times the two
block transports when its indices lie in the same
$L$-block, and zero otherwise. The first two rows
have range below $1$, the third range at most $L-1$.
At rows in $\widetilde U$ their columns and read
bonds stay in its $M$-cube. They agree on both tori.

At zero, $F_0$ has only its diagonal, the four chamber
neighbors, and its block-penalty columns. The hopping
part of its kernel is
\[
 48\Big(
  \bigl(\mu+\tfrac{\rho}{2}\sum_{i=0}^3w_i\bigr)\delta_{x,x'}\mathbf1_4
       -\sum_{j=0}^3\delta_{x',x^{(j)}}M_j(x)U_B(x,x^{(j)})
       \Big).
\]
The factor $48$ is the inverse column weight.
Only the fourth neighbor reads a bond meeting
$\widetilde U$; the penalty reads paths inside the
$L$-block of $c$. All columns lie over
$\widetilde U^{\sim1}$ and its range is at most $L-1$.
Thus this entire row agrees at zero.

These equalities also identify the mass form
with the covering sum:
\[
 \langle\bar\psi,\psi\rangle_{\operatorname{pr}_i\widetilde\square}
 =(\eta_k^3/48)\sum_{c\in\widetilde\square,g,\beta}
                          \bar\psi_\beta(c,g)\psi_\beta(c,g).
\]
At zero the two halves of the first-step vertex have
the common coefficient $1/96$, the same two boundary
transports, and the same generators over the cube
and at its outward pairs. Hence the first vertex
also agrees under $\iota_{\widetilde Y}$.

The remaining local operations can be compared in
this common free algebra. Equality of elements gives
equality of their constant coefficients and of their
antisymmetric quadratic kernels: expand in ordered
covering monomials and use the common measure
$\eta_k^3/48$. For a small polymer, the adjustment
$\Theta_{X,B}$ uses averages over shortest chains
of $M$-cubes in $X$. Common faces, path membership
and transports are identical on its lift, as are
its volume and the coefficients $d_X^\alpha$.
Therefore
\begin{equation}\label{eq:rgcov-adjustment-intertwining}
 \iota_{\widetilde X}\Theta_{X_1,B_1}
                      =\Theta_{X_2,B_2}\iota_{\widetilde X}.
\end{equation}

If $\widetilde U$ is admissible at level $k+1$,
then $L\widetilde U$ is admissible at level $k$.
Its box diameter is $L\operatorname{diam}(\widetilde U)$,
and $N(L\widetilde U)\subset L N(\widetilde U)$.
The scaling rules on the common generators are
$\psi(x)\mapsto L^{-1}\psi(x/L)$ and
$\chi(x,x')\mapsto L^{-1-\alpha}\chi(x/L,x'/L)$.
They preserve the chamber label and the period signs.
Thus
\begin{equation}\label{eq:rgcov-scaling-intertwining}
 S_L\iota_{L\widetilde U}=\iota_{\widetilde U}S_L.
\end{equation}

Finally, the integrations agree on any local element.
For $W$, expand in monomials. A monomial moment is
zero unless its modes are paired; otherwise it is
its ordering sign times $\nu^{-\#\mathrm{pairs}}$.
The mode indices and the weight $\nu$ agree.
At zero this weight is $1/48$, so the factor is
$48^{\#\mathrm{pairs}}$; other modes integrate to $1$.
This proves
\begin{equation}\label{eq:rgcov-W-integration}
 \iota_{\widetilde Y}\int\dd\mu_I(W)\,F
                    =\int\dd\mu_I(W)\,\iota_{\widetilde Y}F.
\end{equation}
The restricted photon law is a product of the same
normalized Gaussian on $[-p_{0,k},p_{0,k}]$ in each
coordinate. Consequently its marginal on common
coordinates is identical. With $d_V=2-2L^{-3}$,
its vacuum contribution on either torus is
\begin{equation}\label{eq:rgcov-box-vacuum}
 \delta\varepsilon_k^0
 =-d_V\log\left((2\pi)^{-1/2}
             \int_{-p_{0,k}}^{p_{0,k}}\exp(-z^2/2)\dd z\right).
\end{equation}
This number depends on $e_k$ and $L$, and is independent
of the side.

\subsection{Extraction and determinant comparison}
\label{subsec:rgcov-determinants}

Suppose $\mathscr E_i\in K_k^x(T_i)$ agree in the sense
of \eqref{eq:rgcov-local-agreement}. For a small
covering polymer $\widetilde X$, their values at $(0,0)$
have the same constant and quadratic coefficients.
Their local extractions
\eqref{eq:rg-local-extractions} are therefore the same:
their integrated kernels and their common volume are
the same. Fixing a covering cube identifies all small
polymers containing its image on either torus. The sums
\eqref{eq:rg-counterterm-maps} consequently give
\begin{equation}\label{eq:rgcov-counterterms}
 m(\mathscr E_1)=m(\mathscr E_2),\qquad
 \varepsilon_{\mathrm v}(\mathscr E_1)
                         =\varepsilon_{\mathrm v}(\mathscr E_2).
\end{equation}
Subtracting the common constant and mass forms preserves
local agreement of $R\mathscr E_i$.
On a small polymer, its quadratic part also agrees;
\eqref{eq:rgcov-adjustment-intertwining} proves agreement
of $(R\mathscr E_i)^{\mathrm{nat}}$.
On other polymers neither of these operations changes
the relevant element.

The determinant terms agree on every flat polymer,
not just on small ones. Let $X_i=\operatorname{pr}_i
\widetilde X$ and suppose $B_i$ have equal lifts on
bonds meeting $\widetilde X$.
For $k\ge1$, set
\[
 O_1=B_{k,r}\gamma_5B_{k,r},\qquad
 O_2=Q_kJ^*S_{k,r}(X_i,B_i)JQ_k^T(-B_i).
\]
The trace in the determinant integral is
\begin{equation}\label{eq:rgcov-local-determinant-trace}
 \operatorname{Tr}(O_1O_2)
       =\sum_{z_1,z_2\in X_i\cap\Z^3}
                      \tr\big(O_1(z_2,z_1)O_2(z_1,z_2)\big).
\end{equation}
Here and below $\operatorname{Tr}$ is the trace on the whole finite
function space, while $\tr$ is the spin trace.
Indeed
\begin{align*}
 O_2(z_1,z_2)
   &=\int_{\Delta_{z_1}}\dd x\int_{\Delta_{z_2}}\dd x'\,
       \varphi_{B_i}(z_1,c)S_{k,r}(X_i,B_i)(x,x')
                            \varphi_{B_i}(z_2,c')^{-1}.
\end{align*}
Local polymer support confines $x,x'$ to $X_i$,
and $\Delta_z$ lies in the $M$-cube of $z$.
Thus both unit sites in this trace lie in $X_i$.
The rows of $B_{k,r}$ stay in their own $L$-block,
so the intermediate sum in $O_1(z_2,z_1)$ also stays
in $X_i$. All sums may now be formed on the lift
$\widetilde X$; the two antiperiodic kernel signs
cancel in each trace summand.
The finite-row comparison and
\eqref{eq:rgcov-local-polymer-kernel} make every factor
the same on both tori.

At zero the required traces are
\begin{align}
 \operatorname{Tr}(F_0(B_i)S_{0,r}(X_i,B_i))
   &=\int_{X_i}\dd x\int_{X_i}\dd x'\,
          \tr\big(F_0(B_i)(x,x')S_{0,r}(X_i,B_i)(x',x)\big),
                                      \label{eq:rgcov-zero-trace}\\
 \operatorname{Tr}(\gamma_5S_{0,r}(X_i,B_i))
   &=\int_{X_i}\dd x\,\tr\big(\gamma_5S_{0,r}(X_i,B_i)(x,x)\big).
                                      \label{eq:rgcov-zero-trace-gamma}
\end{align}
Support again restricts all arguments to $X_i$.
Both are the same covering sums with measure $1/48$
and all $48$ chamber labels. The local $F_0$ and
$S_{0,r}$ kernels agree, so the traces agree for
every $r\ge0$.

These pointwise identities can be integrated in $r$.
The required domination follows
from the primitive estimates. Since $P_B^2=P_B$,
\[
 B_{k,r}\gamma_5B_{k,r}
  =\gamma_5\big((b_k+\mathrm i\gamma_5r)^{-2}(I-P_B)
                +(b_k+b_L+\mathrm i\gamma_5r)^{-2}P_B\big).
\]
In the component matrix norm,
$|(a+\mathrm i\gamma_5r)^{-1}|_c
\le\sqrt2(a^2+r^2)^{-1/2}$ for $a>0$.
Hence the first trace is bounded by a finite
$B_i,X_i$ dependent constant times
$(b_k^2+r^2)^{-1}$.
The other factor is uniformly bounded by the local
resolvent estimate and the two finite block transports.
This majorant is integrable because $b_k>0$.
At zero, $|S_{0,r}(X_i,B_i)(x,x')|_c
\le C_w(1+r)^{-1}$.
The first trace in \eqref{eq:rgcov-zero-trace} is
integrated against $\dd r/r$ on $[1,\infty)$,
and \eqref{eq:rgcov-zero-trace-gamma} against $\dd r$
on $[0,1]$; both are integrable.
The kernels are Borel in $r$, so these are
ordinary absolutely convergent scalar integrals.

The two unsymmetrized determinant primitives agree.
Applying the argument at $-B_i,\bar B_i,-\bar B_i$
identifies all four terms in its symmetrization.
It gives
\begin{equation}\label{eq:rgcov-determinant-agreement}
              E_k^{d,(1)}(X_1,B_1)=E_k^{d,(2)}(X_2,B_2).
\end{equation}
Subtracting the values at $\mathscr C_i$ or at zero
proves the corresponding equality for
$\delta E_k^{\mathrm{det}}$ and $E_k^{\mathrm{det}}$.

\subsection{Comparison of the local fluctuation sums}
\label{subsec:rgcov-fluctuation-comparison}

Let $m$ be real with $|m|\le w_k^{\mathrm m}$, let
$\|\mathscr E_i\|_k^x\le w_k^{\mathrm E}$, and assume
$\mathscr E_i=0$ at zero and local agreement as above.
Fix $A_1\in\widetilde R_k^-(T_1)$ and
$\mathscr C_1\in\mathfrak X_k(T_1)$, and put
$(A_2,\mathscr C_2)=p^*(A_1,\mathscr C_1)$.
All small-charge and primitive hypotheses of the
nonlinear step are imposed on both tori.
In particular the contour radii satisfy
$1<r_2\le r_1$, $r_4>1$, and the domain inclusions
of Proposition~\extref{5.1}{NonlinearStep} hold.
We compare the terms of that step on a connected
nonempty $\widetilde U\subset\widetilde Y$.

Assign common coordinate values
$z_\Upsilon\in[-p_{0,k},p_{0,k}]$ to the photon basis
elements contained in $\widetilde U$, and set all
other coordinate values to zero on each torus.
Assign common parameter values of modulus at most
$M^{\alpha_0}$ to the parameter cubes contained in
$\widetilde U$, and set all other parameters to zero.
These data will be denoted by $\zeta_i,s_i$.
They belong to the respective box and polydisk
because indices on the lift are injective.

The fine fluctuation fields have the same lifts on
$\widetilde U^{\sim(1+d_U)}$:
\begin{equation}\label{eq:rgcov-fine-fluctuation}
                         Z_1(s'_1)(b)=Z_2(s'_2)(b).
\end{equation}
To check the successive supports in their composition,
write it in kernels,
\[
 Z_i(s'_i)(b)
  =\sum_{b',\Upsilon,\Upsilon'}
    H_k^{\mathrm{ph},(i)}(s'_i)(b,b')
          C_i(b',\Upsilon)
            C_k^{1/2,(i)}(s'_i)(\Upsilon,\Upsilon')\zeta_i(\Upsilon').
\]
At such a row $b$, \cref{prop:rgcov-supported-rows}
with $j=1+d_U$ confines $b'$ to
$\widetilde U^{\sim2d_U}$ and identifies the row kernels.
The coordinate row has columns in the $M$-cubes of
the two ends of $b'$, hence in
$\widetilde U^{\sim(1+2d_U)}$; those finite rows agree.
The square-root row, with $j=1+2d_U$, confines
$\Upsilon'$ to $\widetilde U^{\sim3d_U}$.
All these sets lie in $N(\widetilde U)$, and the
last coordinate values agree there.
The two compositions are therefore the same finite
covering sum. At zero the first row is diagonal
and the sum over $b'$ has one term.

Put $B_i=A_i+\mathscr C_i$ and
$B_i(t)=B_i+tZ_i(s'_i)$.
For $|t|\le r_1$ these fields lie in
$\widetilde{\mathcal G}_k(T_i)$, and
$A_i+tZ_i(s'_i)\in\widetilde R_k(T_i)$.
They have equal lifts on the bonds just compared.
The fermion mean-field rows at sites at distance
at most $\eta_k$ from $\widetilde U$ consequently agree,
as do the $\Gamma_k$ rows at unit sites in
$\widetilde U$ and the $M_k(X_i,B_i(t))$ kernels
for lifted subpolymers $\widetilde X\subset\widetilde U$.
Here the parameters vanish outside $U_i$.
The row-vanishing and component argument in the proof
of Proposition~\extref{5.3}{NonlinearStep} confines
their columns to $U_i$ itself: a nonzero mean-field
row has its active parameter component attached to
that set; the $\Gamma_k$ row starting in this component
has the same confined columns.
For $k\ge1$, the parameter-independent term $\mathcal B_{k,0}$
in $\Gamma_k$ also preserves $U_i$, a union of partition
$M$-cubes.
Thus all these fields are polynomials in the common
generators over $\widetilde U$.

For example, at a relevant chamber site,
\[
 (H_k^{(i)}(s_i'',B_i(t))\Psi)_\beta(x)
   =\sum_{y\in\widetilde U\cap\Z^3,\beta'}
         H_k^{(i)}(s_i'',B_i(t))(x,y)_{\beta\beta'}
                                      \Psi(y,\beta',0)
\]
for $k\ge1$, with the analogous barred sum.
Their coefficients agree.
At every pair admissible for $\widetilde X$, applying
the common finite difference row preserves this equality.
The differences between these fields at $t=1$ and $t=0$
are precisely the ordinary and pair components of
$J_k^\#\Psi$ used in $G_2,G_3$.
Both values are in the disk because $r_1>1$.
Moreover locality of $R\mathscr E_i$ gives equality
of its values at the pairs
$(A_i+tZ_i,\mathscr C_i)$ on each subpolymer.
These pairs need only agree on its read bonds,
by the criterion following
\eqref{eq:rgcov-local-agreement}.

It follows that the three first-localization contour
integrands agree under $\iota_{\widetilde U}$.
For $G_1$, this is the common local $M_k$ bilinear.
For $G_2$, addition of the common odd field $tJ_k^\#\Psi$
to each generator commutes with an algebra homomorphism;
apply it to the equal local activity.
For $G_3$, the same substitution applies to the equal
mass form, and $m'=m+m(\mathscr E_i)$ is common by
\eqref{eq:rgcov-counterterms}.
Their contour radii and scalar contour measures
are the same. Integration thus gives
\begin{equation}\label{eq:rgcov-first-G}
              \iota_{\widetilde U}G_a^{(1)}(s_1,X_1)
                                      =G_a^{(2)}(s_2,X_2),
                    \qquad a=1,2,3.
\end{equation}
At zero $G_2=G_3=0$ and no auxiliary field occurs.
$G_1$ is the first-step vertex; its two factors
$\exp(\pm\mathrm i qZ_i(c,c_3))-1$ agree by
\eqref{eq:rgcov-fine-fluctuation}, with its common
coefficient $1/96$ and the common outward pair.
This proves \eqref{eq:rgcov-first-G} at that level too.

The finite localization commutes with this equality.
Indeed the pairs $(X,P)$ with $Y'(X,P)=U_i$ correspond
bijectively to lifted pairs $(\widetilde X,\widetilde P)$
inside $\widetilde U$ having
$Y'(\widetilde X,\widetilde P)=\widetilde U$.
If a torus pair has that union, its polymer and all
its parameter cubes are contained in $U_i$, so the
unique lifts give the inverse correspondence.
The identities for $P_0,Y'$ proved above show that
it preserves the parameter prescriptions.
Its summand is therefore
\begin{equation}\label{eq:rgcov-localization-term}
 \int_{[0,1]^{\widetilde P}}\dd u\,
       \prod_{\mathrm{pc}\in\widetilde P}\partial_{u_{\mathrm{pc}}}
                       G_i(s_i[u],\operatorname{pr}_i\widetilde X),
\end{equation}
where $s_i[u]$ is $1$ on $P_0(\widetilde X)$, $u$ on
$\widetilde P$, and zero elsewhere.
Identity \eqref{eq:rgcov-first-G} holds throughout
the parameter polydisk. Since $M^{\alpha_0}>1$,
the integration cube is in its interior, and
coefficient differentiation and integration preserve
the equality. This proves equality of
$\operatorname{Loc}G_a(U_i)$.

The reconstruction substitution replaces each $\Psi(y)$
by the same two-piece row $T^{(1)}\psi+T^{(2)}\chi$.
Both rows, including their boundary membership indicators,
have already been compared. Their resulting ordinary
generators are over $\widetilde U$, and their boundary
pairs are admissible for it. Equality on generators
proves equality of the substitution homomorphisms.
We obtain
\begin{equation}\label{eq:rgcov-first-localization}
              \iota_{\widetilde U}E_{\le3}^{(1)}(U_1)
                                      =E_{\le3}^{(2)}(U_2).
\end{equation}
At zero no reconstruction substitution is made.

For the second localization, the common fermion field is
\[
 \mathrm{Wc}=H_k(s'',B)\Gamma_k(s'',B)W,\qquad
 \overline{\mathrm{Wc}}=\bar H_k(s'',B)\bar W .
\]
For $k\ge1$ both are sums over the common unit sites
of $\widetilde U$. For $k=0$ they are
\[
 \mathrm{Wc}(x)=d(c)\frac1{48}
            \sum_{x'\text{ over }\widetilde U}
                                      \Gamma_0(s'',B)(x,x')W(x'),
 \qquad
 \overline{\mathrm{Wc}}(x)=d(c)\bar W(x),
\]
where the common diagonal factor $d(c)$ vanishes
outside $U_i$. These formulas preserve all chamber
labels and the weight $1/48$.
The kernel equalities proved above, followed by the
common pair difference rows, identify $W_k^\#$ at
every index used by a subpolymer.

The first-localized activity on a subpolymer reads
only photon coordinates supported there. Hence
\eqref{eq:rgcov-first-localization}, proved with coordinates
supported on that subpolymer, still holds when coordinates
are supported on the larger $\widetilde U$.
Adding the common $W_k^\#$ to its generators proves
equality of $G_{\le3}^{\,W}$.
The same addition to $R\mathscr E_i$ and to the mass
form proves equality of the contour integrands of
$G_4,G_5$. The second localization uses only $\pi_k$;
the correspondence and derivative argument
\eqref{eq:rgcov-localization-term} apply with this
one kind. Consequently
\begin{equation}\label{eq:rgcov-finite-F}
                       \iota_{\widetilde U}F_1(U_1)=F_2(U_2).
\end{equation}

We now pass to the local connected logarithm.
Its finite Mayer activity is
\[
 K_i(X_i)=\sum_{\{\widetilde X_a\}:\,\bigcup_a\widetilde X_a=\widetilde X}
                \prod_a\big(\exp(F_i(\operatorname{pr}_i
                                                    \widetilde X_a))-1\big),
\]
where the $\widetilde X_a$ are distinct connected
subpolymers of $\widetilde X$.
The products are even. Their index sets and factors
are the same on both tori by
\eqref{eq:rgcov-finite-F}; coordinate locality allows
their evaluation with common coordinates on
$\widetilde X$.
For the integrated activity $H_i(X_i)$,
\eqref{eq:rgcov-W-integration} identifies the Wick
integrals. The other photon coordinates integrate
to $1$, leaving the same product of restricted normal
laws over the coordinates supported in
$\widetilde X$. Hence $\iota_{\widetilde X}H_1(X_1)=H_2(X_2)$.

The logarithm has the expansion
\[
 E_k^{\#,(i)}(X_i)=
   \sum_{n\ge1}\frac1{n!}
    \sum_{\bigcup_{a=1}^n\widetilde X_a=\widetilde X}
       u^T(\operatorname{pr}_i\widetilde X_1,\ldots,
                    \operatorname{pr}_i\widetilde X_n)
                      \prod_{a=1}^n H_i(\operatorname{pr}_i\widetilde X_a).
\]
Projection preserves incompatibility inside this lift,
so the graph determining $u^T$ is the same.
Every summand agrees under the common algebra map.
Both series converge absolutely in the weighted norm
by Proposition~\extref{6.1}{NonlinearStep}.
The map $\iota_{\widetilde X}$ preserves that norm,
so equality passes to their sums:
\begin{equation}\label{eq:rgcov-logarithm-agreement}
       \iota_{\widetilde X}E_k^{\#,(1)}(X_1;A_1,\mathscr C_1)
                      =E_k^{\#,(2)}(X_2;A_2,\mathscr C_2).
\end{equation}

\subsection{Agreement of the step after reblocking}
\label{subsec:rgcov-step-comparison}



\begin{proof}[Proof of \Cref{thm:rgcov-step-comparison}]
The scalar conclusions are
\eqref{eq:rgcov-counterterms} and
\eqref{eq:rgcov-box-vacuum}.
For the outgoing activity it suffices to compare a
level $k+1$ connection pair with its pullback.
Let $\widetilde U$ be connected and admissible there,
put $Y_i=\operatorname{pr}_i'\widetilde U$,
and set $\widetilde Y=L\widetilde U$.
The latter is a connected union of $LM$-cubes,
hence also of $M$-cubes, admissible at level $k$.
Its neighborhood is contained in $L N(\widetilde U)$.

If $(\mathscr A',\mathscr C')$ satisfies the outgoing-domain hypotheses at level
$k+1$, its scaled pair lies in
$\widetilde R_k^-\times\mathfrak X_k$ by
\eqref{eq:rg-scaled-domain-inclusion}.
Scaling and pullback commute on covering bonds:
both give $L^{-1/2}B'(b/L)$.
Denote the scaled pairs by $(A_i,\mathscr C_i)$.
They satisfy exactly the hypotheses used in
\eqref{eq:rgcov-logarithm-agreement}.

The polymers with $\bar X^L=LY_i$ are precisely
$\operatorname{pr}_i\widetilde X$, where
$\widetilde X$ is connected in $\widetilde Y$ and
its $LM$-closure is $\widetilde Y$.
Indeed $X\subset\bar X^L$ gives its unique lift
inside $\widetilde Y$. Closure commutes with projection,
and lies in $\widetilde Y^{\sim(L-1)}$.
Projection is injective on that enlargement,
so equality of its projected closure with
$\operatorname{pr}_i\widetilde Y$ forces equality
of the lifted closure with $\widetilde Y$.
This also proves the converse and shows that the
reblocking index set is common.

In that index set the outgoing activity is
\begin{align}
 \mathscr E_i^+(Y_i;\mathscr A_i',\mathscr C_i')
  =\sum_{\substack{\widetilde X\subset\widetilde Y\text{ connected}\\
                         \overline{\widetilde X}^{\,L}=\widetilde Y}}
 S_L\Big[&
     (R\mathscr E_i)^{\mathrm{nat}}
                  (\operatorname{pr}_i\widetilde X;A_i,\mathscr C_i)
       +\delta E_k^{\mathrm{det},(i)}
                  (\operatorname{pr}_i\widetilde X;A_i,\mathscr C_i)
       \nonumber\\
       &+E_k^{\#,(i)}
                  (\operatorname{pr}_i\widetilde X;A_i,\mathscr C_i)\Big].
                                      \label{eq:rgcov-reblocked-comparison}
\end{align}
Each of the three entries agrees under
$\iota_{\widetilde Y}$, by extraction and adjustment,
\eqref{eq:rgcov-determinant-agreement}, and
\eqref{eq:rgcov-logarithm-agreement}, respectively.
The entries belong to its ordinary and pair subalgebra.
Equation~\eqref{eq:rgcov-scaling-intertwining} identifies
their scaled values, so the sums agree under
$\iota_{\widetilde U}$.
Locality of the outgoing family, already proved in
the step theorem, then extends the comparison
from pullback pairs to pairs with equal local lifts.

The external scalar recursion has the same reblocking
index set. Its incoming scalar term agrees by hypothesis,
and its local determinant term agrees by
\eqref{eq:rgcov-determinant-agreement} at
$(\mathscr C_1,\mathscr C_2)$ and at $(0,0)$.
Its two sums therefore agree.

The argument without an external slot uses precisely
these local operations with $B_i=A_i$.
The determinant difference is then
$E_k^{\mathrm{det}}(X_i,A_i)=E_k^d(X_i,A_i)-E_k^d(X_i,0)$.
Extraction, addition of odd fields, finite localization
and integration do not require that slot. Their
same domain and convergence assertions were established
for this case in the nonlinear-step construction.
Thus the three bracketed terms in
\eqref{eq:rgcov-reblocked-comparison} agree in that
case too, and scaling gives the conclusion.
Restriction of a family with a slot to
$\mathscr C=0$ also commutes with every operation
in the displayed formula.

At $k=0$ the incoming activity and its extracted
terms are zero. The comparison is for the local
determinant difference and the first correction,
with chamber generators, traces and Wick moments
of weight $1/48$, as explicitly included above.
\end{proof}


